\section{The next complete matching problem}\label{sec:outlook}
\subsection{Relation to spectral realizations of the zeros}
Corollary~\ref{cor:spectral} places this objective inside a long line of work
on spectral interpretations of the zeros: the Hilbert--P\'olya expectation as
developed by Berry and Keating \cite{BerryKeating} and surveyed by Sierra
\cite{Sierra}, the Hamiltonian proposal of Bender, Brody and M\"uller
\cite{BBM}, and the trace formula of Connes on the adele class space
\cite{ConnesTrace,CCM}. What distinguishes the present objective is not its
conclusion but its class of candidate pairings: protected line and defect
algebras whose positive twisted traces are supplied in advance by a field
theory, rather than an operator constructed to have the required spectrum.
We are not aware of an attempt to realize the Weil form in that class, and
the representation-theoretic results used above
\cite{Sphere,Schur,RG,EKRS,Klyuev} do not appear to have been applied to it.
That is a reason to continue; it is not evidence that the matching exists.

Two features of the existing literature bear directly on
Section~\ref{subsec:interference}. Connes and Consani obtain archimedean Weil
positivity from the trace of the scaling action compressed to the orthogonal
complement of the range of cutoff projections \cite{CCWeil}. That is exactly
a mechanism of the kind Proposition~\ref{prop:no-augment} demands: a negative
contact relative to a positive reference arises as
$\norm{(1-\Pi)v}^2=\norm v^2-\norm{\Pi v}^2$, never as an added channel.
Bombieri's analysis of the quadratic functional \cite{Bombieri} likewise
treats contact, poles and primes together rather than as separate positive
pieces. Both are precedents for the structure that
Theorem~\ref{thm:two-channel} shows to be unavoidable.

\subsection{What the exclusion leaves}\label{subsec:leaves}
\paragraph{The contact must arise by projection.}
A positive norm produces a constant more negative than any of its channels
only through a subtraction that is itself a norm, that is through
$\norm{(1-\Pi)v}^2$ for a projection $\Pi$, or more generally through a
compression. Both the gamma and prime data must then live in one space in
which the projection acts, rather than in separate positive summands. In
particular the next preparation should be specified together with the
operator whose range is removed.

\paragraph{The pole term needs a mechanism.}
$P_L$ is a rank-two Hermitian form with one negative direction, and by
\eqref{eq:explicit-formula} it is the contribution of the two poles of
$\zeta$ at the non-real points $\tau=\pm i/2$. Every pairing used in this
manuscript is positive by construction, and no positive channel supplies a
negative direction. A realization must obtain $P_L$ from a subtraction of
the kind just described, or from a sector in which two states are removed:
a ghost or BRST sector, a boundary condition that is not reflection
positive, or a rank-two defect projection. Locating such a mechanism in a
named protected sector is a better use of the next calculation than a
further prime-channel identity.

\paragraph{What to compute in a candidate model.}
By Corollary~\ref{cor:spectral} a realization compatible in $L$ is
translation covariant with spectral measure the zero counting measure. The
object to extract from a candidate theory is therefore a spectral measure or
density of states, and the question to ask of it is whether it is that
measure. Reproducing the terms of \eqref{eq:target} one at a time is not a
weaker version of the same question: Theorem~\ref{thm:two-channel} shows that
every term can be matched exactly while no realization exists.

\subsection{A finite-prime phase identity is useful background}
For comparison with earlier local-factor approaches, define
$S_p=J_p/\overline{J_p}$ using \eqref{eq:prime-phase}, and set
$S_{\mathcal P}=S_\infty\prod_{p\in\mathcal P}S_p$ for a finite prime
set. Logarithmic differentiation gives
\begin{equation}\label{eq:finite-phase}
 -i\overline{S_{\mathcal P}}S_{\mathcal P}'
 =b(\tau^2)+w_0
 -2\sum_{p\in\mathcal P,\,m\geq1}(\log p)p^{-m/2}
                         \cos(m\tau\log p).
\end{equation}
When $\mathcal P$ contains the active primes of $I_L$, integration
against $|\widehat{E_Lf}|^2/(2\pi)$ and addition of $P_L[f]$ gives
$Q_L[f]$. Inactive repetitions vanish by support. This argument uses
no infinite Euler product on the critical line.

The identity is classical local-factor algebra. Its phase derivative
is negative at $\tau=0$, and the signed poles are outside it. It is
neither a positive norm nor an identification of a canonical semilocal
positive pairing with the Weil form. The present inverse problem is
to find a physical source for the whole form, not merely for this phase.

\subsection{A precise open problem}
\begin{problem}[Joint arithmetic preparation]\label{prob:joint}
Specify a positive quantum, boundary or defect model, with a physically
defined source $\Phi_L(f)$ on smooth compact support, for which the
entire pairing is \eqref{eq:target}. Derive the source's adjoint,
domain and normalization, including the fixed contact, both signed
poles, every active prime repetition, and absence or exact cancellation
of mixed translations. For an all-support realization, give one
compatible rule as $L$ increases.
\end{problem}

A candidate preparation must in addition pass
\eqref{eq:interference-bound} for every decomposition into channels whose
norms it computes separately, and must produce the pole term and the
contact from one mechanism rather than from independent positive pieces.

This is an open problem, not an existence conjecture established for
a named field theory. A worthwhile intermediate theorem could be
conditional on construction of a credible specified theory, provided
the construction assumptions do not prescribe the desired covariance.
Proposition~\ref{prop:conditional} would then give the forward
implication to RH. Neither a reverse implication nor an equivalence
with Yang--Mills is proposed here.

The two leading directions have complementary roles. Sphere boundary
modules give the clearest real Gaussian and oscillator normalization
control. The next sphere object should be a boundary or vortex
correspondence acting on that control with a derived infinite source.
An arbitrary choice of coefficients in an orthogonal basis would not
establish the required physical correspondence. The Schur direction
now has a specified positive prime source, and its next object should
be a charge-neutral gauge or RG/interface preparation.

In the conformal $SU(2)$ benchmark, the magnetic generator has two
shifted contributions and a neutral coefficient, schematically
\begin{equation}\label{eq:neutral-word}
 H=uC_+(v)u+C_0(v)+u^{-1}C_-(v)u^{-1},
\end{equation}
and the representation includes a Weyl constraint
\cite[Sections 5.1.2--5.1.3]{Schur}. The notation in
\eqref{eq:neutral-word} describes the structure to retain; it is not
yet a specified arithmetic source. The exact coefficients and projection
must be carried through a chosen neutral line or interface word. The
neutral term and projection are possible mechanisms for joint
interference, not terms whose desired effect can be assumed. This
calculation must also implement the real-input lift or replace it by
another derived boundary map.

The immediate target is the complete identity at $L=1$:
\begin{equation}\label{eq:first-target}
 \norm{\Phi_1(f)}^2=K[E_1f]+w_0\norm f^2+P_1[f]
 -\frac{2\log2}{\sqrt2}\re\ip{E_1f}{U_{\log2}E_1f}.
\end{equation}
The gamma response, contact, pole terms and prime atom must be computed
together, and by Theorem~\ref{thm:two-channel} they cannot be computed
apart. At $L=5/4$ the atom at $\log(3/2)$ must be absent; at $L=3/2$
the second return of $2$ must have primitive coefficient $(\log2)/2$
for each oriented translation, while the prime $3$ term remains active.
These are complete identity tests. Failure calls for a different source
or model rather than a growing sequence of residual estimates.

\subsection{Scope of the working manuscript}
The positive gamma action, its integer refinement, and the positive
prime reference are established ingredients. The dressed calculation of
Section~\ref{sec:dressed} places the last of these in an explicit elementary
Schur preparation. Section~\ref{subsec:interference} then shows that this
preparation cannot be completed in the way \eqref{eq:full-comparison}
suggests: no coherent source has the gamma energy and the summed prime
references as its two channel norms. The source tests in this manuscript are
sufficiently explicit to be checked and reused, but they have not supplied
the full matching identity. The main remaining difficulty has therefore not
been reduced to a numerical sign check.

Other constructions remain available: interacting boundary gluing,
integrable defects, gauge networks and arithmetic graph models may
provide different source rules. Their broader literature comparison
is retained in the investigation's research notes, while this working
manuscript focuses on the calculations above. The evidence for sphere
and Schur systems is the availability of explicit protected
representations and positive pairings, not an assertion that their
arithmetic matching is inevitable. Further drafts should strengthen a
specific joint construction or replace it when its complete pairing fails.
